\documentclass[10pt]{article}
\usepackage[margin=20mm]{geometry}
\usepackage{amsmath,amssymb,mathtools,fontspec,unicode-math,xurl,hyperref}
\IfFontExistsTF{DejaVu Serif}{\setmainfont{DejaVu Serif}}{\setmainfont{Latin Modern Roman}}
\setmathfont{latinmodern-math.otf}
\setlength{\parindent}{0pt}\setlength{\parskip}{6pt}\sloppy
\title{Technical Review of Polynomial
Moment Inequality Certificates}
\author{AI-assisted technical review}
\date{}
\begin{document}
\maketitle
Abstract: The reviewed 36-page author manuscript contains two false scope statements, an invalid empty-set reduction in its perturbative Positivstellensatz, a non-preserving specialization in a classical corollary, and a false printed covariance-certificate identity. Exact counterexamples identify the failed statements or proof steps. The covariance identity is repaired by a substitution already present in the acquired notebook and verified independently. The bilocal optimum of four has direct analytic and finite support verification. The overall supported assessment is Major, with the two peripheral scope defects assessed individually as Minor. The principal noncompact theorem and its classical corollary are not shown false. Scientific adjudication remains pending.\par
\section*{Source, framework and review scope}
The paper is Igor Klep, Victor Magron and Jurij Volčič, Sums of Squares Certificates for Polynomial Moment Inequalities, DOI https://doi.org/10.1007/s10208-025-09703-z. The reviewed source is the undated 36-page author PDF at https://igorklep.github.io/files/mompoly.pdf; physical and printed page numbers coincide. The 43-page publisher version and 35-page arXiv v2 were identified by metadata, not substituted or certified equivalent.\par
The algebra \(\mathcal M=\mathbb R[\mathrm{m}_\alpha:\alpha\ne0]\) has \(\mathrm{m}_0=1\). The map \(\mathrm{m}:\mathcal M[x]\to\mathcal M\) is \(\mathcal M\)-linear and sends \(x^\alpha\) to \(\mathrm{m}_\alpha\). For \(S_1\subseteq\mathbb R[x]\), let \(K=K(S_1)\) be its nonnegative locus. For \(S_2\subseteq\mathcal M\), \(\mathcal K=\mathcal K(S_1,S_2)\) consists of probability measures supported on \(K\), with all marginal moments finite, satisfying \(s(\mu)\ge0\) for \(s\in S_2\). The pointwise evaluation is \(f(\mu,X)\), where \(X\in K\) is independent of \(\mu\). Random-variable evaluation instead produces a function \(f[\pi,F]\), ordered almost everywhere (pp. 6–8).\par
The cone \(\operatorname{QM}(S_1,S_2)\) consists of finite sums of terms \(h^2s\), \(h^2\mathrm{m}(g^2s)\), and \(h^2t\), with \(s\in\{1\}\cup S_1\), \(t\in S_2\), and moment-polynomial coefficients as in Definition 2.2. Its generators are nonnegative on \(\mathcal K\times K\). No global finiteness assumption on \(S_1\) appears in these definitions; Section 5 explicitly imposes finite constraint sets for computation.\par
The review concentrates on the evaluation bridge, the compact Positivstellensatz and its bounded-set corollary, both Bell applications, and the final noncompact theorem and classical corollary. The Kadison–Dubois and compact moment-representation inputs to Theorem 4.2 were used conditionally, not independently reproved. The pseudo-moment and Hilbert-17th results and the full analytic machinery of Section 6 were not verified.\par
\section*{Problem 1: almost-everywhere evaluation does not imply the stated pointwise property}
Location and assessment: Proposition 2.4, p. 8; Minor. The reverse implication asserts that nonnegativity of \(f[\pi,F]\) for every admissible random variable implies nonnegativity of \(f(\mu,X)\) for every independent \(X\in K\). Taking the coordinate random variable under \(\mu\) only gives nonnegativity \(\mu\)-almost everywhere.\par
Set \(n=1\), \(S_1=\{1-x^2\}\), \(S_2=\{-\mathrm{m}_2\}\), and \(f=-x^2\). Then \(K=[-1,1]\). Every admissible measure satisfies \(\int x^2\,d\mu\le0\), hence is \(\delta_0\). Every admissible random variable has \(E F^2=0\), so \(F=0\) almost everywhere and \(f[\pi,F]=0\) almost everywhere. Thus condition (ii) holds. However, \[f(\delta_0,1)=-1<0,\qquad (\delta_0,1)\in\mathcal K\times K,\] so condition (i) fails. All moments exist, both constraint sets are finite, and even the ordinary quadratic module is archimedean.\par
There is no full-support hypothesis. Adding a null-probability point does not change an almost-everywhere assertion. A precise repair is to replace \(X\in K\) by \(X\in\operatorname{supp}\mu\): for fixed \(\mu\), continuity of \(X\mapsto f(\mu,X)\) makes nonnegativity almost everywhere equivalent to nonnegativity on the support. Pushforward of \(\pi\) proves the corresponding equivalence. The original forward implication also remains valid.\par
This changes the claimed evaluation bridge and the comparison on p. 11. Theorem 4.2 is proved directly for pointwise evaluation, and the Bell objectives are pure moment polynomials, hence constant in the free variable. No failure of those core applications follows from this example; that limited dependency supports Minor.\par
\section*{Problem 2: boundedness alone does not control an infinite presentation}
Location and assessment: Corollary 4.5, p. 16; Minor. The proof infers that the preordering generated by \(S_1\) is archimedean from boundedness of \(K(S_1)\), without assuming \(S_1\) finite.\par
Take \(n=1\), \(S_2=\varnothing\), and \[S_1=\{s_j=x^2(x^2-j):j=1,2,\ldots\}.\] The feasible set is \(K(S_1)=\{0\}\): if \(x\ne0\), choosing an integer \(j>x^2\) makes \(s_j(x)<0\). Thus \(\mathcal K=\{\delta_0\}\), and \(f=-\mathrm{m}_2\) is zero on \(\mathcal K\times K\), satisfying condition (i).\par
Let \(\widetilde S_1\) contain the square-free finite products from \(S_1\). Suppose condition (ii) held with \(\varepsilon=1\), so \(1-\mathrm{m}_2\in\operatorname{QM}(\widetilde S_1,\varnothing)\). Any such identity contains only finitely many indices \(j\), say a finite set \(J\). Choose \(R^2>\max(\{1\}\cup J)\). Evaluate the identity at \((\delta_R,R)\). Every used \(s_j(R)\) is positive; products, squares and their moments are nonnegative. The right side is therefore nonnegative, but the left side is \(1-R^2<0\), a contradiction. The test pair need not satisfy the unused constraints: it tests a claimed finite polynomial identity.\par
Every finite subsystem has an unbounded feasible set, explaining the failure of the asserted archimedean inference. A finite-\(S_1\) interpretation would exclude this witness, but it is not the printed hypothesis. The later explicit discussion of infinitely many constraints on p. 33 reinforces that this is a real scope boundary. A verified modification is to adjoin a valid ball constraint to \(S_1\); the enlarged ordinary module is archimedean and Theorem 4.2 applies to the enlarged cone. This does not establish membership in the original cone. The main archimedean theorem and the finite SDP examples are not refuted, supporting Minor.\par
\section*{Problem 3: the empty-feasible-set reduction is not nonnegative}
Location and assessment: Theorem 6.9, Case (b), p. 31; Major for an invalid derivation of a stated core guarantee. Theorem 6.9, p. 30, permits arbitrary \(S_1\), finite \(S_2\), and \(f\in\mathcal M[x]\), and equates positivity on \(\mathcal K\times K\) with membership of \(f+\varepsilon(\Phi_r+\Psi_r)\) in the moment quadratic module for some \(r\), for every \(\varepsilon>0\). Here \(\Phi_r\) and \(\Psi_r\) are the ordinary and moment perturbations defined in Section 6.\par
For empty \(\mathcal K\), the proof introduces \(z=x_{n+1}\), \(S'_1=z(\{1\}\cup S_1)\) and \(S'_2=\mathrm{m}_{1,0,\ldots,0}^{\,2}S_2\), and asserts \(zf\ge0\) on the new feasible set in order to apply Case (a). The printed moment index refers to the first original coordinate.\par
Take \(n=1\), \(S_1=\varnothing\), \(S_2=\{-1\}\), and \(f=-1\). The original feasible measure set is empty, \(S_2\) is finite, and \(f\) is even an ordinary polynomial, as requested within Case (b). The new sets are \[S'_1=\{z\},\qquad S'_2=\{-\mathrm{m}_{1,0}^{\,2}\}.\] The probability measure \(\nu=\delta_{(0,1)}\) has all moments, is supported where \(z\ge0\), and satisfies \(-(\int x_1\,d\nu)^2=0\). For the admissible point \(X=(0,1)\), nevertheless, \[(zf)(\nu,X)=-1.\] This also fails under random-variable evaluation, so Problem 1 does not explain it away. The necessary premise for invoking the nonempty Case (a) is false.\par
This is a counterexample to the proof step, not to the theorem’s conclusion. Indeed, in this particular original example \(-1\) is already a cone generator. The dependency is the advertised noncompact Positivstellensatz in its empty-set case, not a peripheral notation statement. No complete repair of this reduction is claimed; changing an index alone would require a new argument for all support and moment constraints.\par
\section*{Problem 4: zero-moment specialization does not preserve the cone}
Location and assessment: Corollary 6.14, p. 32; Major for the invalid derivation of the general classical extension. The corollary asserts, for arbitrary \(S\subseteq\mathbb R[x]\), that \(f\ge0\) on \(K(S)\) is equivalent to \(f+\varepsilon\Phi_r\in\operatorname{QM}(S)\) for a suitable \(r\) at every positive \(\varepsilon\). Its proof fixes the free variables and sends every nonconstant moment to zero, claiming \(\zeta(\operatorname{QM}(S,\varnothing))\subseteq\operatorname{QM}(S)\).\par
Take \(S=\{x^2-1\}\) in one variable. By the definition of the moment cone, \(\mathrm{m}(x^2-1)=\mathrm{m}_2-1\) belongs to \(\operatorname{QM}(S,\varnothing)\). Its image is \(-1\). But \(-1\notin\operatorname{QM}(S)\): at the feasible point \(x=1\), every ordinary quadratic-module element is nonnegative. Thus the claimed cone inclusion fails for a finite constraint set with nonempty feasible set. Correcting the proof’s reference from Theorem 4.2 to Theorem 6.9 would not correct this generator failure.\par
For nonempty \(K(S)\), evaluation of moments at a feasible Dirac measure \(\delta_a\), rather than at \(\delta_0\), does preserve the cone: the image of \(\mathrm{m}(g^2s)\) is a nonnegative scalar because \(s(a)\ge0\). Handling the resulting perturbation constants still requires an argument, and the empty-set case requires separate treatment. This report does not certify a full repair. The advertised extension to general classical constrained polynomials is the dependency supporting Major. The conclusion itself has not been disproved.\par
\section*{Problem 5: the printed covariance certificate is a false identity}
Location and assessment: Section 5.1, pp. 19–20; Major for the printed exact proof of the covariance upper bound, with a fully verified algebraic repair. The objective \(f=\operatorname{cov}_{3322}(A,B)\) is the sum of the eight covariance terms on p. 19, with \((x_1,\ldots,x_6)=(A_1,A_2,A_3,B_1,B_2,B_3)\in\{-1,1\}^6\). The displayed \(8\times8\) matrix \(G\) and vector \(v\) on p. 20 are asserted to satisfy \[\mathrm{m}(v^\top Gv)=\tfrac92-f
 \quad\text{modulo }x_j^2=1.\] The first, seventh and eighth entries of \(v\) contain \(\mathrm{m}_{110000}\).\par
Set every random variable equal to \(1\) on a one-point probability space. This is an admissible binary model; Section 5.1 imposes no zero-mean condition. Every moment is one and every covariance zero. The printed vector becomes \[v=(1,0,5,-10,0,1,38/3,38/3)^\top .\] Exact rational multiplication with the printed \(G\) gives \(v^\top Gv=7/2\), whereas \(9/2-f=9/2\). This is a direct failure of the identity, not a violation of the covariance inequality.\par
An independent symbolic expansion, reducing exponents modulo two and applying the formal moment map, gives the discrepancy \[\begin{aligned}
 \mathrm{m}(v^\top Gv)-(\tfrac92-f)
 ={}&(\mathrm{m}_{010000}+\mathrm{m}_{100000}-\mathrm{m}_{110000}) \nonumber\\
 &\cdot(\mathrm{m}_{000001}-\mathrm{m}_{000010}-\mathrm{m}_{000100}
             +2\mathrm{m}_{100000}-2\mathrm{m}_{110000}).\end{aligned}\] The relevant acquired notebook `mompop\_covar\_cert.nb` instead uses \(q=\mathrm{m}_{100000}+\mathrm{m}_{010000}\) in exactly those three positions. Replacing the printed \(\mathrm{m}_{110000}\) by \(q\) and leaving all other entries and \(G\) unchanged makes the symbolic discrepancy identically zero.\par
The leading principal minors of \(G\), calculated exactly, are \[\frac43,\ \frac1{24},\ \frac1{128},\ \frac5{6144},\
 \frac1{12288},\ \frac1{131072},\
 \frac{19}{100663296},\ \frac1{1073741824}.\] They are all positive, proving positive definiteness. The repaired identity therefore establishes \(f\le9/2\) for every binary probability model. The claimed attainment at \(9/2\) in the external reference was not independently reconstructed.\par
This correction is supported by both exact algebra and the archived notebook at the bound commit. No author code was executed, and cached notebook outputs were not treated as new computations. The discrepancy does not establish which code generated the published SDP experiment. It also does not establish author adoption in another manuscript version. The printed derivation is central to settling the stated conjecture, which supports Major even though the exact repair preserves that upper bound.\par
\section*{Positive checks and overall dependencies}
The proof architecture of Theorem 4.2 (pp. 15–16) is consistent within the checked scope: archimedean separation produces a character, and compact moment representation recovers an admissible measure. This uses the paper’s cited external representation theorems conditionally. Problems 1 and 2 do not refute this argument.\par
For the bilocal problem (5.7)–(5.9), pp. 20–22, let \[\eta_1=(1,1,-1,-1),\quad
 \eta_2=(1,-1,1,-1),\quad
 \eta_3=(1,-1,-1,1).\] On uniform independent \(r,s\in\{1,2,3,4\}\), set \(A_i=\eta_i(s)\), \(C_i=\eta_i(r)\), and \(B_i=(1-2\mathbf1_{r=s})\eta_i(r)\). Exact enumeration of all 16 atoms verifies all 64 bilocal factorizations, nine means, twelve off-diagonal pairs, and 21 repeated-index triples. There are zero failed constraints. Moreover \(E B_iC_i=1/2\), \(E A_iB_i=-1/2\), and all six distinct-index triples equal \(-1/2\); the objective is \(4\).\par
The upper bound can also be checked directly. For each \(i\), \[1-E B_iC_i+E A_iB_i
   =E[(1-B_iC_i)(1+A_iB_i)]\ge0,\] because \(E A_iC_i=0\). For distinct \(i,j,k\), write \(a=E A_iA_k\), \(c=E C_iC_k\) and \(T=(A_iA_k+C_iC_k)B_j+A_iC_i+A_kC_k\). Binary identities and the stated vanishing and bilocal constraints give \[E T^2=4+4ac+4U,\qquad E(TB_j)=a+c,
 \quad U=E A_iB_jC_k+E A_kB_jC_i.\] Consequently \(E(T/2-aB_j)^2=1+U\ge0\). Summing these three pair bounds and the three preceding bounds divided by three yields the objective at most \(4\). Thus this Bell conclusion survives the identified defects.\par
A bounded exact-title/correction search and primary metadata lookup on 29 October 2026 located no linked correction. This is neither an exhaustive absence claim nor a first-discovery claim. The publisher lists online publication on 4 March 2025 and pages 1381–1423 in volume 26 (2026); arXiv lists v2 dated 10 May 2024. Their full mathematical text was not compared with this snapshot.\par
Item | Initial-review status\par
Problems 1–2 | Minor; false scope statements with no demonstrated failure of the main archimedean theorem or finite Bell applications.\par
Problems 3–4 | Major; core derivations fail at explicit steps. The theorem and corollary conclusions are not disproved; full repairs remain.\par
Problem 5 | Major printed certificate defect; replacement by the archived first-moment sum verifies the upper-bound repair exactly.\par
Adoption | Unknown; not requested.\par
Remaining work | Full Section 3 and Section 6 verification; empty-set repair; classical corollary proof; publisher-version comparison.\par
Decision | Pending; human reviews empty; initial delivery is not independent acceptance.\par
\section*{Supporting checks}
Manifest SHA256: `b851a857ec523f29520ea95d3b62a564bae925c0f2897465a81f34c7747b9115`. The notebook repository is bound to commit `5c6eb122ea5bfdf05957688ea91f5b52e1d96e3a`. The certificate notebook hash is `73ebf83532f3fa8c1e42e5e8714d09b89237bd9179bab007d8dd17ebc861a968`.\par
All 36 pages were mechanically parsed. Scientific reading covered pp. 1–8, the relevant comparison on p. 11, pp. 15–22, and the theorem/corollary passages on pp. 30–33. The remainder was not scientifically verified. Selected notebook input definitions were read statically; a quoted-token extraction is not a lossless notebook parser, and cached output cells are not runtime evidence.\par
The independent `exact\_checks.py`, run with the existing Python and SymPy, exited zero. Its retained JSON output includes the failed printed identity, the successful repaired identity, all principal minors and all bilocal witness counts. The framework counterexamples are analytic arguments, with their inputs recorded in that JSON; recording them does not convert them into computer proofs. No author code, numerical SDP solver or published experiment was run. The first compilation failed because the cached sans-serif math font metric was unavailable. Its source and full logs are preserved. The notation font was changed to cached serif roman without changing the mathematics. Actual build commands, logs, compiler/dependency payloads and parse results accompany the report. Separate receipts bind the source evidence images and actual viewing of every final report page. This executor performed the scientific work; the helper supplied only clerical checklist advice. The visible product/account observation came from the main dispatcher, while the runtime self-description is recorded separately. Backend equivalence is not certified.\par
9 I. Klep, V. Magron and J. Volčič, Sums of Squares Certificates for Polynomial Moment Inequalities, 36-page author manuscript, exact acquired snapshot, https://igorklep.github.io/files/mompoly.pdf. V. Magron et al., mompop, https://github.com/magronv/mompop, commit `5c6eb122ea5bfdf05957688ea91f5b52e1d96e3a`, selected static notebook inspection only. I. Klep, V. Magron and J. Volčič, same title, Found. Comput. Math. 26 (2026), 1381–1423, DOI: 10.1007/s10208-025-09703-z; metadata consulted. Same authors and title, https://arxiv.org/abs/2306.05761v2, 10 May 2024, 35 pages; metadata consulted, full text not compared.\par
\end{document}
